% Figure 10 draft C, the development drawing: the band with its front third peeled flat to the right in the picture
% plane, the remaining cylinder at the left; then the full chart beneath with the same points and paths.
\documentclass[10pt,border=1mm]{standalone}
\usepackage[T1]{fontenc}
\usepackage{figurestyle}
\input{primitives.tex}
\input{molecules.tex}
\input{numbers.tex}
% the chart of the band: \chartpic{x}{y}{halfwidth cm}{halfheight cm}{fill} draws [-pi,pi] x [-eta0,eta0] with
% the six H-cells, U gold, the six version points, the t and tu paths from H; labels inside the cells.
\newcommand\chartpic[5]{\begin{scope}[shift={(#1,#2)},x=#3,y=#4]
    \path[fill=#5] (-1,-1) rectangle (1,1);
    \foreach \n/\a/\b/\e/\f in \chartCells {\draw[fine] (\a,\e) rectangle (\b,\f);}
    \path[fill=ColRetained!70] (-0.3333,0) rectangle (0.3333,1); \draw[fine] (-0.3333,0) rectangle (0.3333,1);
    \draw[heavy] (-1,-1) rectangle (1,1);
    \node[lab] at (0,.5) {$\mathcal U$}; \node[lab] at (.6667,.5) {$t\,\mathcal U$}; \node[lab] at (-.6667,.5) {$t^2\mathcal U$};
    \node[lab] at (.8333,-.5) {$u\,\mathcal U$}; \node[lab] at (-.8333,-.5) {$u\,\mathcal U$};
    \node[lab] at (-.5,-.5) {$tu\,\mathcal U$}; \node[lab] at (.3333,-.5) {$t^2u\,\mathcal U$};
    \draw[tpath] (0,1)--(.62,1);
    \draw[upath] (0,1)--(-.3333,-1);
    \foreach \x/\y in {0/1,.6667/1,-.6667/1,1/-1,-1/-1,-.3333/-1,.3333/-1} \draw[line width=.5pt,draw=PlateInk,fill=ColDot] (\x,\y) circle[radius=2.1pt];
    \node[sub,anchor=south] at (0,1.06) {$H$}; \node[sub,anchor=south] at (.6667,1.06) {$tH$}; \node[sub,anchor=south] at (-.6667,1.06) {$t^2H$};
    \node[sub,anchor=north] at (-.3333,-1.08) {$tuH$}; \node[sub,anchor=north] at (.3333,-1.08) {$t^2uH$}; \node[sub,anchor=north] at (1,-1.08) {$uH$}; \node[sub,anchor=north] at (-1,-1.08) {$uH$};
    \foreach \x/\l in {-1/{-\pi},0/{0},1/{\pi}} \node[sub,anchor=north] at (\x,-1.3) {$\l$};
    \node[sub,anchor=north] at (.5,-1.3) {$\tau$};
    \node[sub,anchor=east] at (-1.06,1) {$\eta_0$}; \node[sub,anchor=east] at (-1.06,0) {$0$}; \node[sub,anchor=east] at (-1.06,-1) {$-\eta_0$};
    \node[sub,anchor=south,rotate=90] at (-1.42,0) {$\eta$};
  \end{scope}}

\begin{document}
\begin{tikzpicture}[plate]
\path[use as bounding box] (0,0) rectangle (15,11.0);
% the peeled band: cylinder from a=-180 to 180-phi (cut at the front, a=0 is the near point); the arc (0, phi) is
% unrolled into a flat rectangle lying to the right of the cut, in the picture plane
\begin{scope}[shift={(3.6,8.4)}]
  \pgfmathsetmacro{\R}{1.5}\pgfmathsetmacro{\Hh}{.55}\pgfmathsetmacro{\E}{.42}\pgfmathsetmacro{\ph}{120}
  \pgfmathsetmacro{\W}{\R*\ph*pi/180}
  % far half
  \path[fill=ColBand] plot[domain=90:270,samples=40,variable=\a] ({\R*sin(\a)},{\Hh-\E*\R*cos(\a)}) -- plot[domain=270:90,samples=40,variable=\a] ({\R*sin(\a)},{-\Hh-\E*\R*cos(\a)}) -- cycle;
  \foreach \a in {96,104,...,264} {\pgfmathsetmacro{\lw}{0.15+0.35*abs(cos(\a))}\draw[line width=\lw pt,draw=PlateInk!70] ({\R*sin(\a)},{\Hh-\E*\R*cos(\a)})--({\R*sin(\a)},{-\Hh-\E*\R*cos(\a)});}
  \draw[fine] plot[domain=90:270,samples=40,variable=\a] ({\R*sin(\a)},{-\Hh-\E*\R*cos(\a)});
  % near part, from a=-90 to 0 only (the arc 0..phi is peeled away); the generatrices crowd toward the left silhouette
  \path[fill=ColBand] plot[domain=-90:0,samples=30,variable=\a] ({\R*sin(\a)},{\Hh-\E*\R*cos(\a)}) -- plot[domain=0:-90,samples=30,variable=\a] ({\R*sin(\a)},{-\Hh-\E*\R*cos(\a)}) -- cycle;
  \foreach \a in {-88,-82,...,-64} {\pgfmathsetmacro{\lw}{0.2+0.3*((-64-\a)/24)}\draw[line width=\lw pt,draw=PlateInk] ({\R*sin(\a)},{\Hh-\E*\R*cos(\a)})--({\R*sin(\a)},{-\Hh-\E*\R*cos(\a)});}
  \draw[heavy] plot[domain=-90:0,samples=30,variable=\a] ({\R*sin(\a)},{-\Hh-\E*\R*cos(\a)});
  \draw[heavy] plot[domain=-90:0,samples=30,variable=\a] ({\R*sin(\a)},{\Hh-\E*\R*cos(\a)});
  \draw[heavy] (-\R,-\Hh)--(-\R,\Hh);
  \draw[heavy] plot[domain=90:270,samples=40,variable=\a] ({\R*sin(\a)},{\Hh-\E*\R*cos(\a)});
  \draw[heavy] plot[domain=0:90,samples=30,variable=\a] ({\R*sin(\a)},{\Hh-\E*\R*cos(\a)});
  % the cut edge at the front and the peeled flat panel to the right
  \draw[heavy] (0,{-\Hh-\E*\R})--(0,{\Hh-\E*\R});
  \path[fill=ColBand] (0,{-\Hh-\E*\R}) rectangle (\W,{\Hh-\E*\R}); \draw[heavy] (0,{-\Hh-\E*\R}) rectangle (\W,{\Hh-\E*\R});
  \draw[fine] ({\W/2},{-\Hh-\E*\R})--({\W/2},{\Hh-\E*\R});
  % the hidden part of the far wall seen through the gap: the generatrices from a=0 to phi that were peeled
  \draw[hidden] plot[domain=0:\ph,samples=30,variable=\a] ({\R*sin(\a)},{-\Hh-\E*\R*cos(\a)});
  % points: H at the cut (top), tH on the flat panel at a third of a turn, versions in the remaining band
  \hatom{0}{\Hh-\E*\R}{.075}{ColDot} \hatom{\W}{\Hh-\E*\R}{.075}{ColDot} \hatom{\W/2}{-\Hh-\E*\R}{.075}{ColDot}
  \draw[tpath] (0,{\Hh-\E*\R})--({\W*.92},{\Hh-\E*\R});
  \node[sub,anchor=south] at (0,{\Hh-\E*\R+.1}) {$H$}; \node[sub,anchor=south] at (\W,{\Hh-\E*\R+.1}) {$tH$};
  \node[sub,anchor=north] at ({\W/2},{-\Hh-\E*\R-.1}) {$t^2uH$};
\end{scope}
\node[sub,anchor=north,text width=6.4cm,align=center] at (5.0,6.3) {the band of figure 6 cut at $H$ and peeled: a third of a turn lies flat, the rest still curved};
\chartpic{7.5}{2.6}{5.6cm}{1.8cm}{white}
\end{tikzpicture}
\end{document}
